Настоящата дипломна работа разгледа последователно проектирането и разработването на уеб базирана система за електронна търговия с компютърна и офис техника. В рамките на разработката беше създаден реален Java Spring Boot проект, който покрива основните потребителски и бизнес процеси на един компактен онлайн магазин: регистрация, потвърждение на акаунт, вход, поддръжка на профил, управление на фирми, добавяне на продукти, разглеждане на каталог, работа с количка, финализиране на поръчка и поддържане на история на покупки. Една от основните цели на труда беше не просто да се реализира работещо приложение, а да се покаже как подобна система може да бъде структурирана по ясен, анализируем и методологично последователен начин. В този смисъл разработката демонстрира важна инженерна стойност: архитектурата е организирана в отделни слоеве и пакети, домейн моделът е достатъчно ясен и свързан с предметната област, а избраният технологичен стек е едновременно практичен и представителен за съвременна Java уеб разработка.

Проведеният анализ показа, че изборът на Spring Boot, Spring MVC, Thymeleaf, Spring Data JPA и MySQL е напълно оправдан за подобен тип задача. Тези технологии позволяват бързо, но дисциплинирано изграждане на уеб приложение с добра проследимост на логиката от потребителския интерфейс до persistence слоя. Използването на Hibernate Validator, ModelMapper, JavaMail интеграция и \texttt{Pbkdf2PasswordEncoder} допълнително повишава качеството на решението и демонстрира осъзнато използване на наличните средства в екосистемата. От operational гледна точка обаче dependency конфигурацията би могла да бъде по-чиста, особено по отношение на дублираните или твърде свободно декларирани версии, което е една от посоките за бъдеща хигиенизация.

От гледна точка на качествените характеристики проектът се представя на добро ниво за своя мащаб. Архитектурната подреденост е сред най-силните му страни -- налице е ясно разграничение между контролери, service слой, repository интерфейси, entity модели и визуален слой, а използването на отделни binding, service и view модели е свидетелство за осъзнато проектиране, а не за ad hoc имплементация. Поддържаемостта също е на добро ниво: кодът е сравнително четим, класовете имат разпознаваеми роли, а пакетната структура е логична, което позволява на нов разработчик бързо да се ориентира в системата. Използваемостта на приложението е приемлива и дори добра за учебен проект -- основните функции са групирани логично, навигацията е ясно видима и потребителят може да премине през целия жизнен цикъл на покупката без да губи контекст. По отношение на сигурността системата демонстрира няколко правилни решения като хеширане на паролите и забрана за вход на непотвърден потребител, но достъпът се управлява чрез \texttt{HttpSession} и ръчни условни проверки без централизирана security рамка, което я прави достатъчна за учебно-приложна демонстрация, но не и за production внедряване без значително доработване. Разширяемостта е добра, защото добавянето на нова страница следва установения модел: нов контролерен маршрут, нов шаблон, при нужда нов service метод и repository взаимодействие.

Особено значение в рамките на работата има фактът, че проектът не беше разгледан само описателно, а критично. Наред с положителните характеристики бяха идентифицирани и конкретни ограничения, които очертават естествената граница между функционален демонстрационен прототип и production-ready система. Най-значимото от тях е свързано със сигурността: макар паролите да се съхраняват чрез \texttt{Pbkdf2PasswordEncoder}, достъпът до защитените страници се контролира предимно чрез ръчни проверки върху \texttt{HttpSession}, а липсват роля-базирана авторизация, security filter chain и формален security context. Втората важна граница е домейнното моделиране на поръчките -- текущата реализация използва \texttt{HistoryEntity} като обобщение на завършена покупка, но не поддържа отделни позиции, статуси, платежни състояния или жизнен цикъл на поръчката. Същевременно част от логиката за количката и продуктите работи върху глобални колекции, вместо да бъде строго филтрирана по текущия потребител, което в реална многопотребителска среда би довело до логически несъответствия. Наличието на credentials в \texttt{application.properties} е operational компромис, който макар да не пречи на демонстрационното стартиране, ограничава възможността проектът да бъде безопасно споделен или разгръщан без допълнителна подготовка. Front-end организацията остава сравнително фрагментирана с повторение на навигационна лента в множество шаблони, inline CSS и различни версии на външни библиотеки, а локализацията е само частично последователна. Накрая, макар проектът да използва стандартната Maven структура с тестова зависимост, липсва развита стратегия за unit, integration или UI тестове, така че качеството на системата в момента зависи в по-голяма степен от ръчна проверка и кодов анализ.

Тези наблюдения са важни, защото показват, че разработката е осмислена не само като демонстрационен резултат, но и като основа за по-нататъшно професионално развитие. Най-естествената следваща стъпка е пълноценно въвеждане на Spring Security, което би позволило централизиран контрол върху достъпа до маршрути, разграничаване на роли като администратор, оператор и обикновен клиент, използване на стандартизиран authentication flow и по-лесно управление на сесии, logout и защита от неоторизиран достъп. Втората високоприоритетна посока е преработване на модела за количка и поръчки чрез въвеждане на отделни класове от тип \texttt{Order} и \texttt{OrderItem} и по-ясно разграничение между активна количка и историческа поръчка, така че системата да поддържа множество позиции в рамките на една поръчка, статуси, запазване на артикулите в историята без загуба на детайлна информация и по-прецизна отчетност. Operational зрелостта на проекта може значително да се подобри чрез използване на профили за \texttt{development}, \texttt{test} и \texttt{production} среда, изнасяне на чувствителните данни в environment variables или секретно хранилище и въвеждане на миграции на схемата чрез инструмент като Flyway или Liquibase. Front-end слоят би се възползвал от извличане на повторяемите навигационни блокове в Thymeleaf fragments, централизиране на inline CSS във външни стилови файлове, уеднаквяване на използваните версии на Bootstrap, разширяване на локализацията и подобряване на responsive поведението. Накрая, въвеждането на разумна тестова стратегия с unit тестове за service логиката, integration тестове за repository и контролери и сценарийни тестове за регистрация, вход, checkout и история би повишило увереността при бъдещи промени.

Освен приоритетните подобрения съществуват и по-широки направления за развитие, които биха трансформирали системата от учебен прототип към по-зряло приложение: разширяване на каталога с търсене, сортиране, странициране, маркиране на любими продукти и реално качване на изображения; интеграция с платежни оператори, куриерски услуги и по-богат REST API за външни клиенти; както и контейнеризация с Docker и автоматизирано разгръщане чрез CI/CD процес, които биха улеснили повторяемото стартиране на приложението. За по-голяма яснота тези насоки могат да бъдат обобщени в кратка приоритетна пътна карта, представена в таблица \ref{tab:future-roadmap}.

\begin{table}[H]
    \centering
    \caption{Предложена пътна карта за бъдещо развитие}
    \label{tab:future-roadmap}
    \begin{tabularx}{\textwidth}{L{2.5cm}L{4cm}Y}
        \toprule
        Приоритет & Направление & Очакван ефект \\
        \midrule
        Висок & Spring Security и роли & По-сигурен и по-чист модел на достъп \\
        Висок & Рефакториране на количка и поръчки & Коректност в многопотребителска среда и по-богат домейн модел \\
        Висок & Изнасяне на secret конфигурация & По-добра operational безопасност \\
        Среден & Централизация на front-end и локализация & По-добра поддръжка и UX последователност \\
        Среден & Автоматизирани тестове & По-висока надеждност при развитие на кода \\
        Среден & Миграции и профили за среди & По-предсказуем build и deployment процес \\
        Нисък към среден & Разширени търговски функции & По-голяма практическа стойност на приложението \\
        Нисък към среден & Контейнеризация и CI/CD & По-лесно разгръщане и възпроизводимост \\
        \bottomrule
    \end{tabularx}
\end{table}

Развитието на проекта може да бъде планирано стъпково: най-напред следва да се адресират качествата, които пряко влияят на коректността и сигурността, а след това да се премине към по-богати функционални и operational подобрения. От академична гледна точка основните приноси на настоящата работа могат да бъдат обобщени в няколко направления -- систематизирано е реално приложение за електронна търговия в последователна многослойна архитектурна рамка, изведен е ясен модел на предметната област, демонстрирано е как чрез Spring технологии може да бъде реализиран пълен потребителски поток в уеб приложение, извършена е аналитична верификация на функционалността, интерфейса, конфигурацията и operational средата, и са формулирани реалистични насоки за бъдещо развитие към по-сигурна и по-зряла система. Тези приноси потвърждават, че дипломната работа изпълнява своята двойна задача: от една страна представлява практическа инженерна реализация, а от друга -- завършен технически документ, който обосновава използваните решения. Това е особено важно в контекста на специалност като компютърно и софтуерно инженерство, където способността да се премине от идея към добре структурирана реализация и след това към аргументирана документация е съществен показател за професионална зрялост.

Разработеното приложение не следва да бъде интерпретирано като крайна production-ready търговска платформа -- неговата реална стойност е в това, че представлява работещ, концептуално ясен и надграждаем прототип, върху който могат да бъдат наслагвани следващи итерации с по-силна сигурност, по-богат order management модел, по-чист front-end слой, по-добра конфигурационна дисциплина и автоматизирано тестване. Точно тази отвореност към развитие е един от най-ценните резултати на проекта. В заключение може да се обобщи, че поставената цел е постигната: създадена е уеб базирана система за електронна търговия с компютърна и офис техника, която покрива ключов набор от функционалности, реализирана е с подходящи съвременни технологии, описана е в академично издържана форма и е оценена както през призмата на постигнатото, така и през призмата на нейните ограничения. Това прави проекта убедителен завършек на дипломна разработка и стабилна основа за бъдещо инженерно надграждане.
